\begin{abstract}
Outputs on the probability simplex arise as targets (proportions, compositions) and as predictions (classifier softmax vectors), and split conformal prediction offers them marginal coverage for any nonconformity score.
It is less clear how reliably the \emph{choice} of score, such as Euclidean, Hellinger, Aitchison or likelihood-based, can be evaluated, because conditional coverage must be estimated from finite data.
We study this with synthetic Dirichlet and logistic-normal responses, for which the true conditional coverage and the oracle region are known, and with classifier sets on synthetic and CIFAR-10 outputs.
Marginal coverage follows the exact Beta--binomial law for all 15 scores tested and carries no information on which score is better (rank correlation between replication blocks $-0.03$, against $0.94$ for region volume).
The worst-cell coverage estimator is biased low by an amount we give in closed form (mean $-0.026$ for ten cells and 500 test points), the bias differs between scores, and it can reorder scores (Kendall $\tau$ between the mean estimate and the true ranking $0.83$--$0.92$); a single evaluation recovers the true ranking with Kendall $\tau\approx0.5$--$0.7$.
A half-split bias correction, tried because of this, overcorrects and does not improve the ranking.
Normalised and likelihood scores give the smallest regions and flattest conditional coverage on Dirichlet data, whereas Aitchison-distance scores win on homoscedastic logistic-normal data and under-cover near the boundary.
For classifier outputs the same small-cell bias appears; LAC sets are smaller than APS sets (paired interval above zero) in all 36 synthetic cells, and APS has the higher worst size-stratum coverage in 24 of them, the lower in 10.
All results are for known-truth synthetic data, a fixed log-ratio linear mean model and one CNN; we state where they should not be extrapolated.
\end{abstract}
